
\documentclass{article}
\usepackage{colortbl}
\usepackage{makecell}
\usepackage{multirow}
\usepackage{supertabular}

\begin{document}

\newcounter{utterance}

\centering \large Interaction Transcript for game `cladder', experiment `full\_v1.5\_default', episode 6916 with selfrefine{-}Qwen3.5{-}9B.
\vspace{24pt}

{ \footnotesize  \setcounter{utterance}{1}
\setlength{\tabcolsep}{0pt}
\begin{supertabular}{c@{$\;$}|p{.15\linewidth}@{}p{.15\linewidth}p{.15\linewidth}p{.15\linewidth}p{.15\linewidth}p{.15\linewidth}}
    \# & \multicolumn{2}{c}{Player} && \multicolumn{2}{c}{Game Master} \\
    \hline

    \theutterance \stepcounter{utterance}  
    & & & \multicolumn{4}{p{0.6\linewidth}}{
        \cellcolor[rgb]{0.9,0.9,0.9}{
            \makecell[{{p{\linewidth}}}]{
                \texttt{\tiny{[P1$\langle$GM]}}
                \texttt{You are an expert in causal inference. The following question is not a typical commonsense query, but rather a meticulously designed question created by a professor specializing in causal inference, intended to assess the students' mastery of the course content.} \\
\texttt{Imagine a self{-}contained, hypothetical world with only the following conditions, and without any unmentioned factors or causal relationships: Jyka has a direct effect on hwax. Lirg has a direct effect on hwax. Method 1: We look at how jyka correlates with lirg case by case according to hwax. Method 2: We look directly at how jyka correlates with lirg in general. To understand how jyka affects lirg, is it more correct to use the Method 1 than Method 2?} \\
\texttt{Start your answer with 'Yes' or 'No', followed by additional reasoning or evidence to support your explanation.} \\
            }
        }
    }
     \\ \\

    \theutterance \stepcounter{utterance}  
    & & \multicolumn{4}{p{0.6\linewidth}}{
        \cellcolor[rgb]{0.95,0.95,0.95}{
            \makecell[{{p{\linewidth}}}]{
                \texttt{\tiny{[P1$\leftrightarrow$P]}}
                \texttt{{[}DRAFT{]}} \\
\texttt{No.} \\
\\ 
\texttt{In this hypothetical world, **Jyka** and **Lirg** are both direct causes of the common effect **hwax**. This structure creates a **collider** at *hwax* (Jyka \$\textbackslash{}rightarrow\$ hwax \$\textbackslash{}leftarrow\$ Lirg).} \\
\\ 
\texttt{According to the principles of causal inference, specifically the **Collider Bias** (or Berkson's Paradox) described by the "Causal Markov Condition" and d{-}separation:} \\
\texttt{1.  **Method 2 (Marginal Correlation):** In the absence of conditioning on the collider (*hwax*), the path between Jyka and Lirg is blocked. Therefore, Jyka and Lirg are marginally independent, and their correlation should be zero (or reflect only background noise).} \\
\texttt{2.  **Method 1 (Conditional Correlation):** By looking at the correlation "case by case according to hwax," you are conditioning on the collider (*hwax*). Conditioning on a collider **opens** the path between its parents, inducing a spurious statistical association (correlation) between Jyka and Lirg, even though there is no direct causal link between them.} \\
\\ 
\texttt{The question asks if it is more correct to use Method 1 to understand how Jyka **affects** (causally influences) Lirg. Since there is no causal arrow from Jyka to Lirg in this world, the true causal effect is zero. Method 1 produces a non{-}zero correlation due to selection bias, which would lead one to falsely conclude that Jyka affects Lirg. Method 2 correctly reflects the lack of a direct causal relationship. Therefore, Method 1 is incorrect for inferring causality in this scenario.} \\
            }
        }
    }
    & \\ \\

    \theutterance \stepcounter{utterance}  
    & & \multicolumn{4}{p{0.6\linewidth}}{
        \cellcolor[rgb]{0.95,0.95,0.95}{
            \makecell[{{p{\linewidth}}}]{
                \texttt{\tiny{[P1$\leftrightarrow$P]}}
                \texttt{{[}SELF{-}FEEDBACK{]}} \\
\texttt{FORMAT: ok} \\
\texttt{RULES: ok} \\
\texttt{STRATEGY: ok} \\
\texttt{REASON: The answer correctly identifies the collider structure, explains why conditioning on the collider induces spurious correlation, and concludes that Method 1 is incorrect for inferring causality.} \\
            }
        }
    }
    & \\ \\

    \theutterance \stepcounter{utterance}  
    & \multicolumn{4}{p{0.6\linewidth}}{
        \cellcolor[rgb]{0.561,0.737,0.561}{
            \makecell[{{p{\linewidth}}}]{
                \texttt{\tiny{[P1$\rangle$GM]}}
                \texttt{No.} \\
\\ 
\texttt{In this hypothetical world, **Jyka** and **Lirg** are both direct causes of the common effect **hwax**. This structure creates a **collider** at *hwax* (Jyka \$\textbackslash{}rightarrow\$ hwax \$\textbackslash{}leftarrow\$ Lirg).} \\
\\ 
\texttt{According to the principles of causal inference, specifically the **Collider Bias** (or Berkson's Paradox) described by the "Causal Markov Condition" and d{-}separation:} \\
\texttt{1.  **Method 2 (Marginal Correlation):** In the absence of conditioning on the collider (*hwax*), the path between Jyka and Lirg is blocked. Therefore, Jyka and Lirg are marginally independent, and their correlation should be zero (or reflect only background noise).} \\
\texttt{2.  **Method 1 (Conditional Correlation):** By looking at the correlation "case by case according to hwax," you are conditioning on the collider (*hwax*). Conditioning on a collider **opens** the path between its parents, inducing a spurious statistical association (correlation) between Jyka and Lirg, even though there is no direct causal link between them.} \\
\\ 
\texttt{The question asks if it is more correct to use Method 1 to understand how Jyka **affects** (causally influences) Lirg. Since there is no causal arrow from Jyka to Lirg in this world, the true causal effect is zero. Method 1 produces a non{-}zero correlation due to selection bias, which would lead one to falsely conclude that Jyka affects Lirg. Method 2 correctly reflects the lack of a direct causal relationship. Therefore, Method 1 is incorrect for inferring causality in this scenario.} \\
            }
        }
    }
    & & \\ \\

    \theutterance \stepcounter{utterance}  
    & & & \multicolumn{2}{p{0.3\linewidth}}{
        \cellcolor[rgb]{0.9,0.9,0.9}{
            \makecell[{{p{\linewidth}}}]{
                \texttt{\tiny{[GM$|$GM]}}
                \texttt{No} \\
            }
        }
    }
    & & \\ \\

    \theutterance \stepcounter{utterance}  
    & & & \multicolumn{2}{p{0.3\linewidth}}{
        \cellcolor[rgb]{0.9,0.9,0.9}{
            \makecell[{{p{\linewidth}}}]{
                \texttt{\tiny{[GM$|$GM]}}
                \texttt{no} \\
            }
        }
    }
    & & \\ \\

    \theutterance \stepcounter{utterance}  
    & & & \multicolumn{2}{p{0.3\linewidth}}{
        \cellcolor[rgb]{0.9,0.9,0.9}{
            \makecell[{{p{\linewidth}}}]{
                \texttt{\tiny{[GM$|$GM]}}
                \texttt{game\_result = WIN} \\
            }
        }
    }
    & & \\ \\

\end{supertabular}
}

\end{document}
